% !TeX program = lualatex
% Compile from the presentation directory with:
%   lualatex presentation.tex
% Edit the presenter and date below as needed.
\documentclass[aspectratio=169,11pt]{beamer}

\usepackage{fontspec}
\usepackage{luatexja-fontspec}

\setsansfont{Arial}
\setsansjfont{HaranoAjiGothic-Medium}

\renewcommand{\familydefault}{\sfdefault}

\usepackage{amsmath,amssymb}
\usepackage{booktabs}
\usepackage{graphicx}
\usepackage{array}

\usetheme{default}
\setbeamertemplate{navigation symbols}{}
\setbeamercolor{normal text}{fg=black,bg=white}
\setbeamercolor{structure}{fg=black}
\setbeamercolor{alerted text}{fg=black}
\setbeamercolor{title}{fg=black}
\setbeamercolor{subtitle}{fg=black}
\setbeamercolor{author}{fg=black}
\setbeamercolor{date}{fg=black}
\setbeamercolor{itemize item}{fg=black}
\setbeamercolor{itemize subitem}{fg=black}
\definecolor{HeaderBar}{HTML}{213A3D}
\definecolor{PlotBlue}{HTML}{1F77B4}
\definecolor{PlotOrange}{HTML}{FF7F0E}
\definecolor{PlotGreen}{HTML}{CFE8CF}
\definecolor{PlotGreenDark}{HTML}{2E7D32}
\definecolor{PlotRed}{HTML}{D62728}
\setbeamercolor{frametitle}{fg=white,bg=HeaderBar}
\setbeamerfont{frametitle}{series=\bfseries}
\setbeamercolor{footline}{fg=black,bg=white}

\graphicspath{{assets/}}

\title{\textbf{SI Research Discussion}}
\author{志田浩太}
\date{July 27, 2026}

\setbeamertemplate{footline}{
  \leavevmode
  \hbox{%
    \begin{beamercolorbox}[
      wd=\paperwidth,
      ht=2.3ex,
      dp=1.1ex,
      rightskip=0.5cm
    ]{footline}
      \hfill\insertframenumber/\inserttotalframenumber
    \end{beamercolorbox}%
  }
}

\begin{document}

% -----------------------------------------------------------------------------
\begin{frame}[plain,noframenumbering]
  \titlepage
\end{frame}

% -----------------------------------------------------------------------------
\begin{frame}{進捗}
  \begin{itemize}
    \item \text{非線形特徴選択に起因する選択バイアスとランダム選択の検証をした。}
    \vspace{0.55cm}
    \item \text{RF の Tree SHAP で選択された特徴量の選択領域の数値計算・可視化をした。}
    \vspace{0.55cm}
    \item \text{選択的 $p$ 値を求めるコードが書けた。}
  \end{itemize}
\end{frame}

% -----------------------------------------------------------------------------
\begin{frame}{Unadjusted vs. Random (1/4)：実験設定}
  \textbf{グローバル帰無仮説}
  \[
    \begin{gathered}
      X_{ij}\overset{\mathrm{i.i.d.}}{\sim}\mathcal N(0,1),\qquad
      y_i\overset{\mathrm{i.i.d.}}{\sim}\mathcal N(0,1),\qquad
      X\perp y,\qquad \sigma^2=1.
    \end{gathered}
  \]

  \vspace{0.15cm}
  \begin{itemize}
    \item $n=100$, $d=20$。
    \item 選択数 $k=1,\ldots,10$。
    \item 反復数 $R=1000$。
    \item 有意水準 $\alpha=0.05$。
  \end{itemize}

  \vspace{0.18cm}
  \noindent\rule{\linewidth}{0.6pt}
  \vspace{0.14cm}
  \textbf{比較する選択方法}
  \begin{itemize}
    \item \textbf{Random：}$y$ と独立に $k$ 特徴を一様無作為抽出。
    \item \textbf{Unadjusted SHAP：}同じ $y$ で学習した RF の SHAP 重要度上位 $k$ 特徴を選択。
  \end{itemize}
\end{frame}

% -----------------------------------------------------------------------------
\begin{frame}{Unadjusted vs. Random (2/4)：RFとSHAPによる特徴選択}
  \footnotesize
  \setlength{\abovedisplayskip}{5pt}
  \setlength{\belowdisplayskip}{5pt}
  \textbf{選択パイプライン}
  \[
    \begin{gathered}
      (X,y)\longrightarrow
      \widehat f=\operatorname{RF}(X,y)
      \longrightarrow \text{Tree SHAP }\{\phi_{ij}\}_{i,j},\\[0.15cm]
      \{\phi_{ij}\}_{i,j}\longrightarrow
      I_j=\frac1n\sum_{i=1}^{n}|\phi_{ij}|
      \longrightarrow \widehat S=\operatorname{TopK}(I_1,\ldots,I_d).
    \end{gathered}
  \]

  \begin{itemize}
    \item Tree SHAP の $\phi_{ij}$ は観測 $i$ における特徴 $j$ の予測への寄与を表す。
    \item 絶対値の平均により，正負の寄与を相殺せず重要度へ集約する。
  \end{itemize}

  \noindent\rule{\linewidth}{0.6pt}
  \vspace{0.02cm}
  \textbf{Random Forest の三条件}

  \vspace{0.02cm}
  \centering
  {\small
  \renewcommand{\arraystretch}{1.02}
  \begin{tabular}{lccc}
    \toprule
    & Baseline & Larger forest & Regularized forest \\
    \midrule
    Number of trees & 50 & 200 & 100 \\
    Maximum depth & 5 & Unlimited & 3 \\
    Minimum samples per leaf & 1 & 1 & 5 \\
    Feature fraction per split & 1.0 & 1.0 & 0.7 \\
    \bottomrule
  \end{tabular}}
\end{frame}

% -----------------------------------------------------------------------------
\begin{frame}{Unadjusted vs. Random (3/4)：Bスプライン効果の検定}
  \small
  \setlength{\abovedisplayskip}{5pt}
  \setlength{\belowdisplayskip}{5pt}

  \begin{itemize}
    \item 選択特徴 $x_j$ から，中心化3次Bスプラインの正規直交基底 $Q_j\in\mathbb R^{n\times r}$ を作る。
  \end{itemize}

  帰無仮説と検定統計量を次のように定義する：
  \[
    H_{0,j}:Q_j^\top\mu=0,
    \qquad
    T_j=\frac{\lVert Q_j^\top y\rVert_2}{\sigma}
    \overset{H_{0,j}}{\sim}\chi_r.
  \]

  \begin{itemize}
    \item 既知分散を仮定し，$T_j\sim\chi_r$ を用いるカイ検定を行う。
    \item 各反復で選択した $k$ 特徴の棄却割合を求め，$R=1000$ 回で平均。
  \end{itemize}
\end{frame}

% -----------------------------------------------------------------------------
\begin{frame}{Unadjusted vs. Random (4/4)：既知分散での結果}
  \centering
  % The source image contains unknown- and known-variance panels.  Retain the
  % shared y-axis labels from the far left and display only the known-variance
  % panel on the right.  Trim dimensions use the PNG's 300-dpi natural size.
  \includegraphics[
    height=0.58\textheight,
    clip,
    trim=0bp 0bp 883.2bp 19.2bp
  ]{../outputs/unadjusted_vs_random/unadjusted_vs_random_rf_sensitivity.png}%
  \hspace{0.04cm}%
  \includegraphics[
    height=0.58\textheight,
    clip,
    trim=490.8bp 0bp 2.4bp 19.2bp
  ]{../outputs/unadjusted_vs_random/unadjusted_vs_random_rf_sensitivity.png}

  \vspace{0.02cm}
  \begin{itemize}\small
    \setlength{\itemsep}{0.06cm}
    \item FPRは上位 $k$ 特徴の平均である。
    \item $k$ を増やすと，偶然の関連が比較的弱い下位特徴も平均に加わるため，FPRが低下する。
  \end{itemize}
\end{frame}

% -----------------------------------------------------------------------------
\begin{frame}{選択的推論：1次元経路の構築}
  \small

  4ページ目で定義した検定統計量 $T_j$ を選択後も用いるが，
  SHAPで $j$ が選ばれた事実に条件づけて帰無分布を補正する。
  Bスプライン効果空間への射影を $P_j=Q_jQ_j^\top$ とし，
  \[
    \begin{gathered}
      a_j=(I-P_j)y,\qquad
      u_j=\frac{P_jy}{\lVert P_jy\rVert_2},\\
      T_{\mathrm{obs}}=\frac{\lVert P_jy\rVert_2}{\sigma},\qquad
      y_j(z)=a_j+\sigma u_jz.
    \end{gathered}
  \]

  $a_j$ と方向 $u_j$ を固定し，効果の大きさだけを $z\ge0$ で動かす。

  \vspace{1em}

  $H_{0,j}:Q_j^\top\mu=0$ の下で $a_j$ と方向 $u_j$ を固定すると，
  残る確率変動は半径
  \[
    Z=\frac{\lVert Q_j^\top y\rVert_2}{\sigma}\sim\chi_r
  \]
  のみになる。
\end{frame}

% -----------------------------------------------------------------------------
\begin{frame}{選択的推論：$z$ ごとに特徴選択をやり直す}
  \small

  ここまでで作った1次元経路 $y_j(z)$ に沿って応答を動かし，
  各 $z$ で同じSHAP選択を再現する。

  \[
    y_j(z)
    \ \longrightarrow\ \text{RF 再学習}
    \ \longrightarrow\ \text{Tree SHAP}
    \ \longrightarrow\ \widehat S(z).
  \]

  $A_j(z)$ は，各 $z$ で選択条件が成立すれば1，不成立なら0となる関数である。
  選択領域を
  \[
    \mathcal Z_j=\{z\ge0:A_j(z)=1\}
  \]
  と定める。

  \begin{itemize}
    \item 各 $z$ で RF の学習から top-$k$ 選択までをすべてやり直す。
    \item 選択と検定に同じ $y$ を使った影響を，$A_j(z)=1$ への条件づけで補正する。
  \end{itemize}

  選択後の基準分布は，$\chi_r$ を $\mathcal Z_j$ 上で切断・再正規化した分布になる。
\end{frame}

% -----------------------------------------------------------------------------
\begin{frame}{選択事象：何に条件づけるか}
  \small
  \setlength{\abovedisplayskip}{3pt}
  \setlength{\belowdisplayskip}{3pt}
  検定対象 $j\in\widehat S_{\mathrm{obs}}$ に対して
  $A_j(z)=1$ を定める条件が3通り考えられる。

  例として，観測された上位2特徴を
  $\widehat S_{\mathrm{obs}}=(x_1,x_2)$，検定対象を $j=x_1$ とする。

  \vspace{0.02cm}
  \begin{itemize}
    \setlength{\itemsep}{0.18cm}
    \item \textbf{対象特徴が選ばれる}
      \quad $j\in\widehat S(z)$

      $j$ が上位 $k$ に残ればよい
      （例：$\widehat S(z)=(x_1,x_3)$ なら成立）。

    \item \textbf{上位 $k$ 特徴の集合が一致する}
      \quad $\operatorname{set}(\widehat S(z))
        =\operatorname{set}(\widehat S_{\mathrm{obs}})$

      集合内の順位は問わない
      （例：$\widehat S(z)=(x_2,x_1)$ でも成立）。

    \item \textbf{上位 $k$ 特徴とその順位が一致する}
      \quad $\widehat S(z)=\widehat S_{\mathrm{obs}}$

      特徴と順位をともに固定する
      （例：$\widehat S(z)=(x_1,x_2)$ のときだけ成立）。
  \end{itemize}

  \vspace{1mm}
  \textbf{条件が厳しいほど，対応する選択領域は小さい：}
  \[
    \mathcal Z_j^{(\text{特徴・順位})}
    \ \subseteq\
    \mathcal Z_j^{(\text{特徴集合})}
    \ \subseteq\
    \mathcal Z_j^{(\text{対象特徴})}.
  \]
  \begin{itemize}
    \setlength{\topsep}{2pt}
    \item $k=1$ では，3つの条件は同一。
  \end{itemize}
  \vspace{-3pt}
\end{frame}

% -----------------------------------------------------------------------------
\begin{frame}{選択領域の数値実験：ベースライン設定}
  \small
  \textbf{データ生成}
  \[
    X_{ij},y_i\overset{\mathrm{i.i.d.}}{\sim}\mathcal N(0,1),\qquad
    X\perp y,\qquad n=100,\quad d=20,\quad\sigma=1.
  \]
  seed $1,\ldots,6$ の各データについて， $(I_1,\ldots,I_d)$ の
  上位 $k$ 特徴を選択し，$k=1$ と $k=2$ を比較する。

  \vspace{0.18cm}
  \textbf{Random Forest の設定}
  \begin{center}
    \renewcommand{\arraystretch}{1.12}
    \begin{tabular}{lcccc}
      \toprule
      & Number of trees & Maximum depth & Minimum leaf size & Random seed \\
      \midrule
      Baseline & 50 & 5 & 1 & 42 \\
      \bottomrule
    \end{tabular}
  \end{center}

  \vspace{0.08cm}
  \textbf{選択領域の探索}
  \begin{itemize}
    \setlength{\itemsep}{0.08cm}
    \item 上位 $k$ 特徴の集合の一致（exact-set）に条件づける。
    \item $z\in[0,6.335]$ を1001点で走査し，変化点を二分法で精密化する。
  \end{itemize}
\end{frame}

% -----------------------------------------------------------------------------
\begin{frame}{選択領域の代表例：$k=1$}
  \begin{columns}[T,onlytextwidth]
    \begin{column}{0.49\textwidth}
      \centering
      \includegraphics[
        width=0.96\linewidth,
        height=0.50\textheight,
        keepaspectratio,
        clip,
        trim=0bp 29.6bp 859.2bp 324bp
      ]{../outputs/selection_region_sweep/k1_baseline/selection_regions.png}
    \end{column}
    \begin{column}{0.49\textwidth}
      \centering
      \includegraphics[
        width=0.96\linewidth,
        height=0.50\textheight,
        keepaspectratio,
        clip,
        trim=429.6bp 29.6bp 429.6bp 324bp
      ]{../outputs/selection_region_sweep/k1_baseline/selection_regions.png}
    \end{column}
  \end{columns}

  \vspace{0.04cm}
  \centering
  {\footnotesize
    \fcolorbox{PlotGreenDark}{PlotGreen}{\rule{0pt}{0.11cm}\hspace{0.27cm}}
    $\mathcal Z_j$\quad
    \textcolor{gray}{\rule{0.35cm}{1.4pt}} $\chi_3$\quad
    \textcolor{PlotBlue}{\rule{0.35cm}{1.4pt}} 条件付き密度\quad
    \textcolor{PlotOrange}{\rule{0.35cm}{1.4pt}} $q_{\mathrm{final}}$\quad
    \textcolor{PlotRed}{\rule{0.35cm}{1.4pt}} $T_{\mathrm{obs}}$
  }

  \vspace{0.08cm}
  \begin{itemize}\small
    \item Data set 4：$x_{12}$，$T_{\mathrm{obs}}=1.808$，
      $P(Z\in\mathcal Z_{12})=0.7334$。
    \item Data set 5：$x_9$，$T_{\mathrm{obs}}=3.362$，
      $P(Z\in\mathcal Z_9)=0.0798$。
    \item 同じ設定でも選択領域の大きさはデータごとに大きく異なる。
  \end{itemize}
\end{frame}

% -----------------------------------------------------------------------------
\begin{frame}{選択領域の代表例：$k=2$}
  \centering
  \includegraphics[
    width=0.96\textwidth,
    height=0.52\textheight,
    keepaspectratio,
    clip,
    trim=429.6bp 36.4bp 0bp 980bp
  ]{../outputs/selection_region_sweep/k2_baseline/selection_regions.png}

  \vspace{0.04cm}
  {\footnotesize
    \fcolorbox{PlotGreenDark}{PlotGreen}{\rule{0pt}{0.11cm}\hspace{0.27cm}}
    $\mathcal Z_j$\quad
    \textcolor{gray}{\rule{0.35cm}{1.4pt}} $\chi_3$\quad
    \textcolor{PlotBlue}{\rule{0.35cm}{1.4pt}} 条件付き密度\quad
    \textcolor{PlotOrange}{\rule{0.35cm}{1.4pt}} $q_{\mathrm{final}}$\quad
    \textcolor{PlotRed}{\rule{0.35cm}{1.4pt}} $T_{\mathrm{obs}}$
  }

  \vspace{0.06cm}
  \begin{itemize}\small
    \item Data set 6 の上位2特徴は $(x_7,x_5)$。
      $P(Z\in\mathcal Z_7)=0.8772$，$P(Z\in\mathcal Z_5)=0.2617$。
    \item 同じ選択集合に条件づけても，各特徴に固有の経路
      $y_j(z)$ と選択領域 $\mathcal Z_j$ が得られる。
    \item したがって，選択的 $p$ 値も特徴ごとに計算する。
  \end{itemize}
\end{frame}

% -----------------------------------------------------------------------------
\begin{frame}{AIS (1/4)：選択的 $p$ 値と単純 Monte Carlo の問題}
  \small
  選択的 $p$ 値は，$A_j(Z)=1$ に条件づけた上側確率：
  \[
    \begin{aligned}
      p_{\mathrm{sel}}
      &=\frac{\displaystyle
        \int_0^\infty A_j(z)\mathbf1\{z\ge T_{\mathrm{obs}}\}
        f_{\chi_r}(z)\,dz}
      {\displaystyle
        \int_0^\infty A_j(z)f_{\chi_r}(z)\,dz}\\
      &=\Pr\!\left(Z\ge T_{\mathrm{obs}}\mid A_j(Z)=1\right).
    \end{aligned}
  \]

  \textbf{単純 Monte Carlo の問題}
  \begin{itemize}
    \setlength{\itemsep}{0.18cm}
    \item $\chi_r$ 分布から候補値 $z$ を生成しても，$A_j(z)=0$ が多い。
    \item 各候補で RF 再学習と SHAP 再計算が必要なため，計算が無駄になりやすい。
  \end{itemize}

  \vfill
  \centering
  \textbf{AISでは，選択されやすい領域へ候補を集中させる。}
\end{frame}

% -----------------------------------------------------------------------------
\begin{frame}{AIS (2/4)：pilot による $q_{\mathrm{adapt}}$ の構築}
  \small
  \textbf{記号}
  \begin{itemize}
    \setlength{\itemsep}{0.06cm}
    \item $q_t(z)=q_{\mathrm{TN}}(z;\mu_t,s_t)$：
      $N(\mu_t,s_t^2)$ を $z\ge0$ に切断した，候補値 $z$ を生成するための分布。
    \item $\mu_t,\,s_t$：$t$ 回目の生成分布における平均と標準偏差。
    \item $z_h\sim q_t$：生成分布 $q_t$ から得た $h$ 番目の候補値。
  \end{itemize}

  \begin{enumerate}
    \setlength{\itemsep}{0.08cm}
    \item \textbf{初期化：}
      $\mu_0=T_{\mathrm{obs}}$，
      $s_0=\max(1,\,0.25T_{\mathrm{obs}}+0.5)$。
    \item \textbf{候補生成：}
      $z_h\sim q_t$ ごとに特徴選択をやり直し，$A_j(z_h)=1$ を残す。
    \item \textbf{重み付けと推定：}
      \begin{itemize}
        \setlength{\itemsep}{0.02cm}
        \item 重み：$w_h=f_{\chi_r}(z_h)/q_t(z_h)$。
        \item 重みの役割：$q_t$ から候補を生成したことによる偏りを補正。
        \item 重み付き平均・標準偏差を，次の
          $(\mu_{t+1},s_{t+1})$ の更新に用いる。
      \end{itemize}
    \item \textbf{$q_t$ の更新：}
      \begin{itemize}
        \setlength{\itemsep}{0.02cm}
        \item 選択候補あり：次の $(\mu_{t+1},s_{t+1})$ を，
              現在値 $(\mu_t,s_t)$ から推定値へ半分だけ近づける。
        \item 選択候補なし：平均は変えず，$s_{t+1}=1.5s_t$。
        \item 固定回数だけ更新し，最後の $q_t$ を $q_{\mathrm{adapt}}$ とする。
      \end{itemize}
  \end{enumerate}
\end{frame}

% -----------------------------------------------------------------------------
\begin{frame}{AIS (3/4)：3つの生成分布を混ぜる}
  \footnotesize
  \setlength{\abovedisplayskip}{3pt}
  \setlength{\belowdisplayskip}{3pt}
  最終標本は，次の分布から生成する：
  \[
    q_{\mathrm{final}}(z)
    =0.25f_{\chi_r}(z)
     +0.375q_{\mathrm{obs}}(z)
     +0.375q_{\mathrm{adapt}}(z).
  \]

  \begin{itemize}
    \setlength{\itemsep}{0.10cm}
    \item $f_{\chi_r}$：元の帰無分布 $\chi_r$ の PDF。
    \item $q_{\mathrm{obs}}$：$T_{\mathrm{obs}}$ を中心とする切断正規分布
      \[
        q_{\mathrm{obs}}(z)=q_{\mathrm{TN}}(z;T_{\mathrm{obs}},s_{\mathrm{obs}}),
        \qquad
        s_{\mathrm{obs}}=\max(0.75,\,0.25T_{\mathrm{obs}}+0.5).
      \]
      $T_{\mathrm{obs}}$ 付近と，それより大きい候補を増やす。
    \item $q_{\mathrm{adapt}}$：pilot で見つけた，
      特徴 $j$ が選択されやすい範囲から候補を増やす切断正規分布。
  \end{itemize}

  \textbf{3つを混ぜる理由}
  \begin{itemize}
    \setlength{\itemsep}{0.08cm}
    \item $q_{\mathrm{obs}}$ だけでは，$T_{\mathrm{obs}}$ から離れた選択範囲を見落とし得る。
    \item $q_{\mathrm{adapt}}$ だけでは，pilot が見つけなかった範囲を探索できない。
    \item $f_{\chi_r}$ を25\%混ぜることで，広い範囲を探索する。
    \item $q_{\mathrm{final}}(z)\ge0.25f_{\chi_r}(z)$ となるため，
      重要度重みを $w_i\le4$ に抑え，
      少数の候補に推定結果が左右されるのを防ぐ。
  \end{itemize}
\end{frame}

% -----------------------------------------------------------------------------
\begin{frame}{AIS (4/4)：重み付き推定と信頼性の確認}
  \footnotesize
  \setlength{\abovedisplayskip}{3pt}
  \setlength{\belowdisplayskip}{3pt}
  固定した $q_{\mathrm{final}}$ から最終標本を生成し，
  \[
    z_i\overset{\mathrm{i.i.d.}}{\sim}q_{\mathrm{final}},
    \qquad
    w_i=\frac{f_{\chi_r}(z_i)}{q_{\mathrm{final}}(z_i)}
  \]
  とする。重み $w_i$ で提案分布の変更を補正すると，
  \[
    \widehat p_{\mathrm{sel}}
    =\frac{\sum_iw_iA_j(z_i)
      \mathbf1\{z_i\ge T_{\mathrm{obs}}\}}
    {\sum_iw_iA_j(z_i)}.
  \]

  \noindent\rule{\linewidth}{0.6pt}
  \vspace{-0.02cm}
  \textbf{ESS（有効サンプルサイズ）を確認する目的：}
  少数の極端に大きい重みが $\widehat p_{\mathrm{sel}}$ を支配し，
  推定値が不安定になっていないかを診断する。
  \[
    \operatorname{ESS}(w)=\frac{(\sum_i w_i)^2}{\sum_i w_i^2}.
  \]
  「等重みの標本何個分の情報か」を表す。
  （例：100候補の重みがほぼ10候補に集中すれば
  $\operatorname{ESS}\approx10$，すなわち実質10候補分である。）
  \begin{itemize}
    \setlength{\topsep}{0.02cm}
    \setlength{\itemsep}{0.02cm}
    \item \textbf{分母 ESS：}$A_j(z_i)=1$ の候補の重みから，
      $\widehat p_{\mathrm{sel}}$ の分母を支える情報量を診断する。
    \item \textbf{分子ESS：}さらに $z_i\ge T_{\mathrm{obs}}$
      の候補に絞り，$\widehat p_{\mathrm{sel}}$ の
      分子を支える情報量を診断する。
    \item \textbf{実装設定：}分母 ESS $\ge80$，分子 ESS $\ge15$。
      未達なら最終標本を80個ずつ追加し，最終標本が合計800個でも
      未達なら $p$ 値を返さない。
  \end{itemize}
\end{frame}

% -----------------------------------------------------------------------------
\begin{frame}{選択特徴が複数あるとき：多重検定をどう扱うか}
  \begin{itemize}
    \setlength{\itemsep}{0.25cm}
    \item 上位 $k$ 個の特徴を選択した場合は，各特徴
      $j\in\widehat S_{\mathrm{obs}}$ に固有の経路 $y_j(z)$ 上で
      同じ計算を行い，$k$ 個の選択的 $p$ 値を得る。
    \item 検定対象が事前に定めた1特徴だけなら補正は不要である。
    \item 複数特徴のいずれかを有意と報告する際，
      1件以上の偽陽性が生じる確率（FWER）を有意水準 $\alpha$ 以下に
      抑えたい場合は，Bonferroni 補正などを適用する。
  \end{itemize}
\end{frame}

% -----------------------------------------------------------------------------
\begin{frame}[t,plain,noframenumbering]{Appendix：ベースライン実験の結果一覧}
  \small
  \centering
  \renewcommand{\arraystretch}{1.18}
  \begin{tabular}{cc cc cc cc}
    \toprule
    & & \multicolumn{2}{c}{$k=1$}
      & \multicolumn{2}{c}{$k=2$：第1特徴}
      & \multicolumn{2}{c}{$k=2$：第2特徴} \\
    \cmidrule(lr){3-4}\cmidrule(lr){5-6}\cmidrule(lr){7-8}
    Data & Seed
      & 特徴 & $P(\mathcal Z_j)$
      & 特徴 & $P(\mathcal Z_j)$
      & 特徴 & $P(\mathcal Z_j)$ \\
    \midrule
    1 & 1 & $x_7$  & 0.6505 & $x_7$  & 0.8424 & $x_8$  & 1.0000 \\
    2 & 2 & $x_7$  & 0.4108 & $x_7$  & 0.6099 & $x_{13}$ & 0.2211 \\
    3 & 3 & $x_{10}$ & 0.2689 & $x_{10}$ & 0.3719 & $x_{13}$ & 0.1498 \\
    4 & 4 & $x_{12}$ & 0.7334 & $x_{12}$ & 0.9420 & $x_2$  & 0.1635 \\
    5 & 5 & $x_9$  & 0.0798 & $x_9$  & 0.5388 & $x_4$  & 0.1879 \\
    6 & 6 & $x_7$  & 0.0732 & $x_7$  & 0.8772 & $x_5$  & 0.2617 \\
    \bottomrule
  \end{tabular}

  \vspace{0.40cm}
  \begin{itemize}
    \item $P(\mathcal Z_j)=P(Z\in\mathcal Z_j)$，$Z\sim\chi_3$。
    \item $k=2$ では，同じデータでも選択位置と対象特徴により
      選択領域の確率が大きく異なる。
    \item 全例で $T_{\mathrm{obs}}\in\mathcal Z_j$。
  \end{itemize}
\end{frame}

% -----------------------------------------------------------------------------
\begin{frame}[t,plain,noframenumbering]{Appendix：1次元経路の導出 (1/3)}
  \small
  \setlength{\abovedisplayskip}{5pt}
  \setlength{\belowdisplayskip}{5pt}

  \textbf{Step 1：応答を効果空間とその直交補空間に分ける}

  $y\sim\mathcal N(\mu,\sigma^2I_n)$ とし，$Q_j^\top Q_j=I_r$，
  $P_j=Q_jQ_j^\top$ とおく。$P_j$ は直交射影なので，
  \[
    y=(I-P_j)y+P_jy
      =a_j+P_jy,\qquad a_j=(I-P_j)y
  \]
  と直交分解できる。

  2つの成分の共分散は
  \[
    \operatorname{Cov}(a_j,P_jy)
    =\sigma^2(I-P_j)P_j=0.
  \]
  両者は同時正規なので，無相関なら独立である：
  \[
    \boxed{a_j\ \mathrel{\perp\!\!\!\perp}\ P_jy}
  \]
  したがって，$a_j$ を固定しても $P_jy$ の確率変動は残る。
\end{frame}

% -----------------------------------------------------------------------------
\begin{frame}[t,plain,noframenumbering]{Appendix：1次元経路の導出 (2/3)}
  \small
  \setlength{\abovedisplayskip}{5pt}
  \setlength{\belowdisplayskip}{5pt}

  \textbf{Step 2：効果空間内の変動を「方向」と「大きさ」に分ける}

  帰無仮説 $H_{0,j}:Q_j^\top\mu=0$ の下では
  \[
    g_j=\frac{Q_j^\top y}{\sigma}
    \sim\mathcal N_r(0,I_r).
  \]
  これを方向 $v_j$ と半径 $Z$ に分ける：
  \[
    v_j=\frac{g_j}{\lVert g_j\rVert_2},\qquad
    Z=\lVert g_j\rVert_2\sim\chi_r.
  \]
  等方性により，方向 $v_j$ と半径 $Z$ は独立である。
  $Q_j$ は長さを保つので，$u_j=Q_jv_j$ とおけば
  \[
    P_jy=Q_jQ_j^\top y
         =\sigma ZQ_jv_j
         =\sigma Zu_j,\qquad
    u_j=\frac{P_jy}{\lVert P_jy\rVert_2}.
  \]
  ここで，$u_j$ は効果の方向，$Z$ は効果の大きさである。

  Step 1の独立性も合わせると，$a_j,u_j$ を固定した後に残るのは
  $Z\sim\chi_r$ だけである。したがって，
  \[
    \boxed{y_j(z)=a_j+\sigma u_jz,\qquad z\ge0}
  \]
  という1次元経路を得る。
\end{frame}

% -----------------------------------------------------------------------------
\begin{frame}[t,plain,noframenumbering]{Appendix：1次元経路の導出 (3/3)}
  \small
  \setlength{\abovedisplayskip}{5pt}
  \setlength{\belowdisplayskip}{5pt}

  \textbf{Step 3：観測データが経路上にあることを確認する}

  $Q_j$ は正規直交基底なので，
  \[
    \begin{aligned}
      \lVert P_jy\rVert_2^2
      &=(Q_j^\top y)^\top Q_j^\top Q_j(Q_j^\top y) \\
      &=\lVert Q_j^\top y\rVert_2^2 .
    \end{aligned}
  \]
  よって，観測された半径は
  \[
    T_{\mathrm{obs}}
    =\frac{\lVert P_jy\rVert_2}{\sigma}
    =\frac{\lVert Q_j^\top y\rVert_2}{\sigma}.
  \]
  $z=T_{\mathrm{obs}}$ を経路に代入すると，
  \[
    \begin{aligned}
      y_j(T_{\mathrm{obs}})
      &=a_j+\sigma u_jT_{\mathrm{obs}} \\
      &=(I-P_j)y
        +\sigma\frac{P_jy}{\lVert P_jy\rVert_2}
        \frac{\lVert P_jy\rVert_2}{\sigma} \\
      &=(I-P_j)y+P_jy=y.
    \end{aligned}
  \]

  \vspace{0.15cm}
  \centering
  \textbf{$z=T_{\mathrm{obs}}$ が観測データに対応し，$z$ は効果の大きさだけを動かす。}
\end{frame}

% -----------------------------------------------------------------------------
\begin{frame}[t,plain,noframenumbering]{Appendix：pilot proposal の更新式}
  \small
  \setlength{\abovedisplayskip}{4pt}
  \setlength{\belowdisplayskip}{4pt}

  \textbf{初期値}\quad
  \[
    \mu_0=T_{\mathrm{obs}},\qquad
    s_0=\max\!\left(1,\,0.25T_{\mathrm{obs}}+0.5\right)
  \]

  \textbf{選択された候補から目標位置と広がりを推定}\quad
  $q_t(\cdot)=q_{\mathrm{TN}}(\cdot;\mu_t,s_t)$ から生成した候補値のうち，
  $A_j(z_h)=1$ を満たす添字集合を $\mathcal H_t$ とする。
  ここで，$f_{\chi_r}(z)$ はカイ分布の PDF（確率密度関数），
  $q_t(z)$ は候補を生成する分布の PDF である。
  \[
    w_h=\frac{f_{\chi_r}(z_h)}{q_t(z_h)},\qquad
    \widehat\mu_t
      =\frac{\sum_{h\in\mathcal H_t}w_hz_h}
             {\sum_{h\in\mathcal H_t}w_h},
  \]
  \[
    \widehat v_t
      =\frac{\sum_{h\in\mathcal H_t}w_h
        (z_h-\widehat\mu_t)^2}
             {\sum_{h\in\mathcal H_t}w_h}.
  \]

  \textbf{現在値から推定値へ半分だけ近づける}
  \[
    \boxed{
    \mu_{t+1}=0.5\mu_t+0.5\widehat\mu_t,\qquad
    s_{t+1}
      =\max\!\left(0.25,\,
        0.5s_t+0.5\sqrt{\max(\widehat v_t,0)}\right)}
  \]
  選択候補または有限な重みがない場合は，中心を保ち探索幅を広げる：
  \[
    \mu_{t+1}=\mu_t,\qquad s_{t+1}=1.5s_t
  \]
\end{frame}

% -----------------------------------------------------------------------------
\begin{frame}[t,plain,noframenumbering]{Appendix：積分から $\widehat p_{\mathrm{sel}}$ への導出}
  \small
  \setlength{\abovedisplayskip}{5pt}
  \setlength{\belowdisplayskip}{5pt}

  上側かつ選択されることを
  \[
    B_j(z)=A_j(z)\mathbf1\{z\ge T_{\mathrm{obs}}\}
  \]
  と書く。求めたい選択的 $p$ 値は，2つの積分の比である：
  \[
    p_{\mathrm{sel}}=\frac{N}{D},\qquad
    N=\int_0^\infty B_j(z)f_{\chi_r}(z)\,dz,\qquad
    D=\int_0^\infty A_j(z)f_{\chi_r}(z)\,dz.
  \]

  $q(z)=q_{\mathrm{final}}(z)$，$U\sim q$，
  $w(z)=f_{\chi_r}(z)/q(z)$ とおくと，
  \[
    N=\mathbb E_q[w(U)B_j(U)],\qquad
    D=\mathbb E_q[w(U)A_j(U)].
  \]
  $q$ から生成した $n$ 個の候補値 $z_1,\ldots,z_n$ で期待値を
  標本平均に置き換える：
  \[
    \widehat N=\frac1n\sum_{i=1}^n w_iB_j(z_i),\qquad
    \widehat D=\frac1n\sum_{i=1}^n w_iA_j(z_i).
  \]
  よって $1/n$ が約分され，
  \[
    \boxed{
    \widehat p_{\mathrm{sel}}
      =\frac{\widehat N}{\widehat D}
      =\frac{\sum_iw_iA_j(z_i)
        \mathbf1\{z_i\ge T_{\mathrm{obs}}\}}
             {\sum_iw_iA_j(z_i)}}
  \]
  となる。帽子は，有限個の候補による Monte Carlo 推定値であることを表す。
\end{frame}

\end{document}
